\section{Introduction}

LLM training is, at its core, a pursuit of \textbf{convergence speed} grounded in the necessity of \textbf{stability}. While the community has explored various optimization strategies, we argue that the essential principle governing training stability is the \emph{Maximal Update Parametrization} ($\boldsymbol{\mu}$P)~\citep{yang2023spectral}. By mandating that the spectral norms of weights and updates scale as $\Theta(\sqrt{d_{\text{out}}/d_{\text{in}}})$ to ensure width-invariant activations remains $\Theta(1)$ scale, $\boldsymbol{\mu}$P serves as the mathematical safeguard against activation explosions~\citep{takase2025spikemorestabilizingpretraining}.  However, the current landscape is saturated with methods that fail to satisfy these fundamental conditions. Conventional soft regularization methods, such as decoupled weight decay or initialization strategies, prove insufficient over long training horizons~\citep{kosson2025weightdecaymattermup}. This unconstrained weight drift destabilizes the \emph{effective step size} (update-to-weight ratio) and degrades feature learning.

On the other side of the spectrum lies the pursuit of optimal convergence. The recent Muon optimizer~\citep{jordan2024muon} has demonstrated remarkable efficiency, often interpreted as steepest descent under the spectral norm. In analyzing Muon, we uncover a surprising insight: it acts as a \textbf{``half-aligned''}  solution to the $\boldsymbol{\mu}$P constraints. However, maintaining stable features requires constraining not only the updates but also the weights themselves. Unstable activations like attention logits explosion were still observed in Muon training~\citep{kimiteam2025kimik2openagentic}. Consequently, practitioners are forced to rely on ``non-essential'' architectural patches to artificially force stability~---~ranging from aggressive normalization schemes like SandwichNorm~\citep{ding2021cogviewmasteringtexttoimagegeneration} and QK-Norm~\citep{henry2020querykeynormalizationtransformers}, to ad-hoc fixes like logit softcapping~\citep{kimiteam2025kimik2openagentic}~---~often at the cost of model expressivity and requiring extensive hyperparameter tuning. This observation motivates a fundamental question:
\begin{quote}
\emph{What if an optimizer could simultaneously satisfy the steepest descent property for \textbf{convergence speed} and the strict $\boldsymbol{\mu}$P constraints for \textbf{fundamental stability}?}
\end{quote}
To answer this, we propose the mathematically unique solution that unifies these two objectives. By identifying the spectral sphere as the natural manifold for stable feature learning, SSO derives the steepest descent direction constrained within this geometry. Unlike heuristic manifold projection methods~\citep{xie2025mhcmanifoldconstrainedhyperconnections,pethick2025trainingdeeplearningmodels}, SSO solves a constrained optimization problem in the tangent space via a Lagrange multiplier search, followed by a retraction step to map the trajectory back onto the spectral sphere.

To enable large-scale training, we offer a systematic implementation in Megatron. We provide principled guidelines for spectral preconditioned optimization, specifically deriving the optimal \emph{learning rate scaler}, determining the critical \emph{atomic module granularity}, and identifying the optimal \emph{spectral radius} to control activation at optimal scales precisely. These offer a robust recipe for large-scale training. Specifically to mitigate the overhead of the iterative root solver, we utilize a novel distributed strategy centered on atomic module sharding\citep{emergeing_optimizer}. This technique partitions fused params into independent spectral units, enabling communication-free local updates. We further address solver-induced workload imbalance through a size-aware ping-pong placement strategy, and accelerate matrix operations using adaptive kernel dispatcher, alongside multi-stream execution and singular vector caching.

Empirically, we validate SSO through extensive pretraining experiments across various scales, including Dense 1.7B, MoE 8B-A1B, and 200-layer DeepNet models. SSO consistently outperforms AdamW and Muon while uniquely preserving stable $\boldsymbol{\mu}$P learning rate transfer. Notably, it yields superior training dynamics: it significantly improves MoE router load balancing, suppresses outliers in deep networks, and strictly bounds activations within a tunable scale.
